\section{NNLO double-real sector}
\label{sec:nnlo-double-real}

The NNLO double-real contribution contains two independent phase-space
momenta after momentum conservation.  The amplitude algebra and IBP reduction
are complete for the saved unpolarized process.  The remaining physics problem
is the analytic evaluation of the resulting cut masters and their assembly
with the exact rational coefficients.

\subsection{Master and boundary classification}

Exact IBP reduction gives
\begin{equation}
 N_{\rm master}^{\rm NNLO}=342.
 \label{eq:nnlo-master-count}
\end{equation}
Applying the relations defined in
Section~\ref{sec:analytic-masters} gives
\begin{equation}
 N_{\rm pow}^{\rm NNLO}=130,
 \qquad
 N_{\rm den}^{\rm NNLO}=82.
 \label{eq:nnlo-equivalence-counts}
\end{equation}
One denominator class is the elementary cut-volume class.  Among the other
81 classes, 17 are maximal under inclusion of their positive-power
denominator sets:
\begin{equation}
 N_{\rm max,nontrivial}^{\rm NNLO}=17.
 \label{eq:nnlo-maximal-class-count}
\end{equation}
These 17 classes are useful starting points for differential systems, but
Eq.~\eqref{eq:nnlo-maximal-class-count} is not the number of scalar boundary
constants.  That number is obtained only after the differential systems have
been constructed, their local modes constrained by the physical contour, and
period identities under $\sim_\partial$ established.

\subsection{Exact reconstruction of a hard master coefficient}

The IBP coefficient of one representative NNLO master receives 1129 rational
contributions from the ordered diagram pairs.  After removing an exact common
monomial, the coefficient was reconstructed with a \Ratracer\ operation graph
and \FireFly\ finite-field interpolation.  Reconstructions with independent
prime sequences gave the same rational function.  Evaluating the original
operation graph and the reconstructed result in a further prime field gave
zero difference, and an independent characteristic-zero substitution gave
the same result.  This establishes that the reconstruction is an exact
rational identity for this coefficient.  It does not evaluate the associated
master integral and is not used as evidence that every NNLO coefficient has
already been reconstructed.

\subsection{A coupled eight-master three-particle-cut family}

We now give a fully analytic example with two dimensionless invariants.  Let
the massless momenta obey
\begin{equation}
 k_a^2=k_b^2=k_c^2=0,
 \qquad
 2k_a\cdot k_b=s,
 \qquad
 -2k_a\cdot k_c=t,
 \qquad
 -2k_b\cdot k_c=u,
 \label{eq:nnlo-worked-kinematics}
\end{equation}
and define
\begin{equation}
 P=k_a+k_b-k_c,
 \qquad
 x=-\frac{t}{s},
 \qquad
 y=-\frac{u}{s},
 \qquad
 \Delta=1-x-y=\frac{P^2}{s}.
 \label{eq:nnlo-worked-variables}
\end{equation}
The solution is constructed in the chamber
\begin{equation}
 0<y<(1-x)^2<1-x<1.
 \label{eq:nnlo-worked-chamber}
\end{equation}

The top master is the physical integral
\begin{align}
 G_7(s,t,u;\eps)={}&
 \int\dd^Dk_e\,\dd^Dk_f\,
 \frac{
  \delta_+(k_e^2)\delta_+(k_f^2)
  \delta_+((P-k_e-k_f)^2)
 }{
  (P-k_e)^2(P-k_f)^2
  (k_a-k_c-k_f)^2(k_a-k_e)^2}.
 \label{eq:nnlo-top-physical-integral}
\end{align}
All four ordinary denominators have a fixed sign on the closed physical phase
space.  The first two are invariant masses of sums of future-directed emitted
momenta and are nonnegative.  The last denominator is
$(k_a-k_e)^2=-2k_a\cdot k_e\leq0$.  Momentum conservation writes the third
as $(k_b-k_e-k_g)^2\leq0$, where $k_g=P-k_e-k_f$; the inequality is the
subset case of Appendix~\ref{app:fixed-sign}.  Hence no ordinary $i0$
choice enters Eq.~\eqref{eq:nnlo-top-physical-integral}.

Embed the integral in the family
\begin{align}
 D_1&=k_e^2,&
 D_2&=k_f^2,&
 D_3&=(P-k_e-k_f)^2,\notag\\
 D_4&=k_a\cdot k_f,&
 D_5&=k_b\cdot k_e,&
 D_6&=(P-k_e)^2,\notag\\
 D_7&=(P-k_f)^2,&
 D_8&=(k_a-k_c-k_f)^2,&
 D_9&=(k_a-k_e)^2.
 \label{eq:nnlo-family-denominators}
\end{align}
The first three slots are cuts and $D_4,D_5$ are irreducible scalar
products.  With
\begin{equation}
 I_{\boldsymbol\nu}
 =\int\dd^Dk_e\,\dd^Dk_f
 \prod_{j=1}^{3}\mathcal C_+^{(\nu_j)}(D_j)
 \prod_{j=4}^{9}D_j^{-\nu_j},
 \label{eq:nnlo-worked-family}
\end{equation}
cut-aware IBP closes on
\begin{equation}
\begin{aligned}
 G_1&=I_{111000000},&
 G_2&=I_{111000010},&
 G_3&=I_{111000011},&
 G_4&=I_{112000011},\\
 G_5&=I_{111001010},&
 G_6&=I_{111001011},&
 G_7&=I_{111001111},&
 G_8&=I_{112001111}.
\end{aligned}
\label{eq:nnlo-eight-master-basis}
\end{equation}
The top sector $(G_7,G_8)$ is genuinely coupled.  The second integral has a
differentiated third cut, not a fourth physical cut line.

\subsection{Canonical differential system}

After the homogeneous powers of (s) are removed, exact invariant
differentiation and IBP reduction give
\begin{equation}
 \dd\boldsymbol J
 =\left[A_x(x,y,\eps)\dd x+A_y(x,y,\eps)\dd y\right]
 \boldsymbol J,
 \label{eq:nnlo-raw-differential-system}
\end{equation}
with
\begin{equation}
 \partial_xA_y-\partial_yA_x+A_yA_x-A_xA_y=0.
 \label{eq:nnlo-flatness}
\end{equation}
A rank-one double pole at $\Delta=0$ occurs only in the top block.  It is
removed by
\begin{equation}
 M_{\rm top}
 =\mathbf1_2+
 \frac{2\eps y(1-y)(2+3\eps)}{(1+\eps)\Delta}E_{21}.
 \label{eq:nnlo-moser-transformation}
\end{equation}
The subsequent rational basis change $\boldsymbol J=T\boldsymbol F$ gives
\begin{equation}
 \dd\boldsymbol F
 =\eps\sum_{a=1}^{6}R_a\,\dd\log\phi_a\,\boldsymbol F,
 \qquad
 \{\phi_a\}=
 \{x,1-x,y,1-y,1-x-y,(1-x)^2-y\},
 \label{eq:nnlo-canonical-system}
\end{equation}
where the $R_a$ are constant matrices.  Every letter has fixed sign along
the real transport path used below.

\subsection{Boundary modes and the independent period}

The boundary $y\to0^+$ is nonuniform.  At $x\to0^+$ its physical
Frobenius expansion contains
\begin{equation}
 x^0,
 \qquad x^{-\eps},
 \qquad x^{-2\eps}.
 \label{eq:nnlo-boundary-modes}
\end{equation}
The lower-sector normalizations reduce to beta and gamma functions.  If
\begin{equation}
 \boldsymbol B_\lambda(x)
 =x^\lambda(b_0+xb_1+\mathcal O(x^2)),
 \qquad
 A(x)=\frac{A_{-1}}{x}+A_0+\mathcal O(x),
 \label{eq:nnlo-frobenius-expansion}
\end{equation}
then
\begin{equation}
 \left[(\lambda+1)\mathbf1-A_{-1}\right]b_1=A_0b_0.
 \label{eq:nnlo-frobenius-recurrence}
\end{equation}
The boundary-to-canonical map has
$C(x)=C_{-1}/x+C_0+\mathcal O(x)$, so the finite vector is
\begin{equation}
 f_\lambda=C_0b_0+C_{-1}b_1.
 \label{eq:nnlo-finite-boundary-match}
\end{equation}
Keeping only $b_0$ would miss a finite contribution.

After the lower modes are fixed, one independent top-sector period remains.
A manifestly real representation is
\begin{align}
 r_{\rm top}(\eps)={}&-
 \frac{\Gamma(2-2\eps)\Gamma(3-3\eps)}
      {\eps\Gamma(-2\eps)\Gamma(1-\eps)^3}
 \int_0^1\dd\alpha\,\dd\beta\,\dd\gamma\;
 \alpha^{-1-\eps}\beta^{-1-\eps}\gamma^{-1-\eps}
 \notag\\[-1mm]
 &\times(1-\alpha)^{-1-2\eps}(1-\beta)^{-1-2\eps}
 (1-\alpha\beta)^{2\eps}(1-\alpha\beta\gamma)^\eps.
 \label{eq:nnlo-corner-period}
\end{align}
Every powered base is positive for $0<\alpha,\beta,\gamma<1$.  Exact
endpoint resolution and hyperlogarithmic integration give
\begin{align}
 r_{\rm top}(\eps)={}&-\frac{10}{\eps^3}
 +\frac{65}{\eps^2}-\frac{135}{\eps}+90-8\zeta_3
 \notag\\
 &+\eps\left(52\zeta_3-\frac{2\pi^4}{3}\right)
 +\eps^2\left(\frac{13\pi^4}{3}-108\zeta_3-288\zeta_5\right)
 +\mathcal O(\eps^3).
 \label{eq:nnlo-corner-period-series}
\end{align}
This is the only new period in this family beyond the beta- and gamma-function
data.

\subsection{Analytic solution and numerical comparison}

The canonical boundary is transported along
\begin{equation}
 (0,0)\longrightarrow(x,0)\longrightarrow(x,y).
 \label{eq:nnlo-transport-path}
\end{equation}
The first segment has letters $\{0,1\}$; the second has
$\{0,1,1-x,(1-x)^2\}$.  The chamber
\eqref{eq:nnlo-worked-chamber} places the endpoint below every nonzero letter,
so no singular letter is crossed.  Ordered Chen integrals solve the system,
\begin{equation}
 \boldsymbol F(x,y)
 =U_y(x;y)U_x(x)\boldsymbol F_\partial,
 \qquad
 U=P\exp\!\left[\eps\int\sum_aR_a\dd\log\phi_a\right].
 \label{eq:nnlo-ordered-transport}
\end{equation}
The residue matrices do not commute, so the path and word order are part of
the result.

Normalize by
\begin{equation}
 \mathcal N_0(\eps)=s^{-1+2\eps}G_1(s,0,0;\eps),
 \qquad
 \widehat G_7=\frac{s^{3+2\eps}G_7}{\mathcal N_0}.
 \label{eq:nnlo-top-normalization}
\end{equation}
Then
\begin{equation}
 \widehat G_7(x,y;\eps)
 =\sum_{n=-3}^{2}\eps^ng_n(x,y)+\mathcal O(\eps^3),
 \qquad g_{-4}=0,
 \label{eq:nnlo-top-series}
\end{equation}
with
\begin{equation}
 g_{-3}=-\frac{5}{\Delta^2},
 \label{eq:nnlo-gminus3}
\end{equation}
\begin{equation}
 g_{-2}=\frac{1}{\Delta^2}
 \left[\frac{65}{2}-4\mathcal G(0;x)+10\mathcal G(1;x)
 -6\mathcal G(1;y)+20\mathcal G(1-x;y)\right].
 \label{eq:nnlo-gminus2}
\end{equation}
Here
\begin{equation}
 \mathcal G(a_1,\ldots,a_n;z)
 =\int_0^z\frac{\dd t}{t-a_1}\mathcal G(a_2,\ldots,a_n;t),
 \qquad \mathcal G(;z)=1,
 \label{eq:nnlo-gpl-definition}
\end{equation}
with tangential regularization at zero.  The coefficients through
$\eps^2$ are exact iterated-integral expressions retained in the ancillary
file \texttt{ancillary/NNLO\_three\_particle\_cut\_master.wl}.

At $s=10$, $x=3/10$, and $y=1/5$, define
\begin{equation}
 R(\eps)=s^4\frac{G_7}{G_1}
 =\sum_{n=-4}^{2}r_n\eps^n+\mathcal O(\eps^3).
 \label{eq:nnlo-fixed-point-ratio}
\end{equation}
The analytic expression and an independently generated \AMFlow\ calculation
give the coefficients in Table~\ref{tab:nnlo-amflow}.
\begin{table}[htbp]
\centering
\small
\begin{tabular}{@{}cr@{}}
\toprule
$n$ & $r_n$\\
\midrule
$-4$ & $0$\\
$-3$ & $-40$\\
$-2$ & $282.32024127181498559344$\\
$-1$ & $-662.90998166471340573018$\\
$0$ & $418.83703822298626042003$\\
$1$ & $-12.21808964614470122794$\\
$2$ & $-32.81177029956454278367$\\
\bottomrule
\end{tabular}
\caption{Laurent coefficients of Eq.~\eqref{eq:nnlo-fixed-point-ratio}.
The analytic and \AMFlow\ values agree to relative precision $10^{-12}$.}
\label{tab:nnlo-amflow}
\end{table}

Two additional stored families have been solved through $\eps^2$ by the
same differential-equation and boundary-period procedure and compared with
independent \AMFlow\ values.  They demonstrate reuse of the method across
distinct families but do not complete the 342-master basis.

\subsection{Remaining analytic work}

The current counts identify which integrals occur and which denominator
classes are maximal.  They do not yet determine all scalar boundary periods
for the 342-master system.  The next analytic step is to construct the
differential closures of the 17 maximal nontrivial classes, determine their
physical local modes, and quotient their boundary data by exact period
identities.  Only after all required masters are known analytically can their
rational coefficients be assembled and the NNLO endpoint distributions be
extracted.
